\begin{frame}
	\frametitle{BitNet b1.58}
	\framesubtitle{Comparado a LLaMA em FP16, mesmo tamanho de modelo \cite{ma2024era}}
	\begin{center}
		\begin{tabular}{lcc}
			\toprule
			& \textbf{3B parâmetros} & \textbf{70B parâmetros} \\
			\midrule
			Memória de GPU & $3{,}55\times$ menor & $7{,}2\times$ menor \\
			Velocidade de decodificação & $2{,}71\times$ mais rápido & $4{,}1\times$ mais rápido \\
			\textit{Throughput} & --- & $8{,}9\times$ maior \\
			Perplexidade / tarefas finais & \multicolumn{2}{c}{equivalente ao FP16} \\
			\bottomrule
		\end{tabular}
	\end{center}
	\vspace{2pt}
	\begin{itemize}
		\item Os autores reportam paridade estatística com a linha de base FP16 a partir de $\sim$3B parâmetros.\newline
		\item \alert{Abaixo dessa escala, a paridade some}: em modelos pequenos, a capacidade perdida pela quantização ternária pesa mais do que o modelo consegue compensar.\newline
		\item De onde vem o nome ``1,58 bit'': $\log_2(3) \approx 1{,}58$ --- o número de bits necessário para representar 3 estados ($-1,0,1$) por peso.
	\end{itemize}
\end{frame}
